\slajdid{09_2-s046}
\begin{frame}[label=09_2-s046]
\color{DEcrvena}
\[\frac{W_p}{W_n}=\frac{\mu_n}{\mu_p}\]
Ako je N tranzistor minimalne geometrije smatra se da je normalizovano $W_n=1$\par\bigskip
Zbog tehnoloških parametara normalizovano $W_p$ mora biti ceo broj, eventualno X.5
\[\frac{W_p}{W_n}=\frac{\mu_n}{\mu_p}\approx\frac{\textit{ceo broj P}}{\textit{ceo broj N}}\]
Često pitanje je kolika je potreba širina $W_p$ da bi prag odlučivanja bio na polovini napona napajanja za poznato $W_n$ u tehnologiji npr 10nm. U tom slučaju $W_n=$ (ceo broj N) x L
\[W_p=\textit{ceo broj}\left(\frac{\mu_n}{\mu_p}\times(\textit{ceo broj N})\right)\times L\]
\end{frame}
